%macbook 
\setcounter{paragraph}{0}
\PhanI
\loai{Điều kiện giao thoa}
\begin{ex}
Hiện tượng giao thoa ánh sáng chỉ quan sát được khi hai nguồn ánh sáng là hai nguồn
\choice
{đơn sắc}
{\True kết hợp}
{cùng màu sắc}
{cùng cường độ}
\loigiai{}
\end{ex}
\begin{ex}%[3P2Y5-1][Loai1] 
Hiện tượng giao thoa sóng xảy ra khi có sự gặp nhau của hai sóng
\choice
{xuất phát từ hai nguồn bất kì}
{xuất phát từ hai nguồn truyền ngược chiều nhau}
{xuất phát từ hai nguồn dao động cùng biên độ}
{\True xuất phát từ hai nguồn sóng kết hợp cùng phương}
\loigiai{
Điều kiện giao thoa: Hai nguồn sóng phải là hai nguồn kết hợp:
\begin{itemize}
\item Cùng phương.
\item Cùng tần số.
\item Hiệu số pha không đổi theo thời gian.
\end{itemize}}
\end{ex} \haidong{20}
\begin{ex}%[3P2Y5-1][Loai1]
\nguon{QG - 2017} Hai nguồn kết hợp là hai nguồn dao động cùng phương, cùng
\choice
{\True tần số và có hiệu số pha không đổi theo thời gian}
{biên độ nhưng khác tần số}
{pha ban đầu nhưng khác tần số}
{biên độ và hiệu số pha thay đổi theo thời gian}
\loigiai{
Hai nguồn kết hợp là hai nguồn có:
\begin{itemize}
\item Cùng phương.
\item Cùng tần số.
\item Hiệu số pha luôn không đổi theo thời gian.
\end{itemize}}
\end{ex} \haidong{20}
\begin{ex}%[3P2Y5-1][Loai1]
Điều kiện để hai sóng cơ khi gặp nhau, giao thoa được với nhau là hai sóng phải xuất phát từ hai nguồn dao động
\choice
{cùng biên độ và có hiệu số pha không đổi theo thời gian}
{cùng tần số, cùng phương}
{có cùng pha ban đầu và cùng biên độ}
{\True cùng tần số, cùng phương và có hiệu số pha không đổi theo thời gian}
\loigiai{
Điều kiện để hai sóng cơ khi gặp nhau, giao thoa được với nhau là hai sóng phải xuất phát từ hai nguồn dao động cùng tần số, cùng phương và có hiệu số pha không đổi theo thời gian.
}
\end{ex} \haidong{20}
\loai{Vị trí cực đại}
\begin{ex}%[3P2Y5-1][Loai1] 
Trong sự giao thoa của hai sóng cơ phát ra từ hai nguồn điểm kết hợp, cùng pha, những điểm dao động với biên độ cực đại có hiệu khoảng cách $d_2 - d_1$ tới hai nguồn, thỏa mãn điều kiện nào sau đây (với k là số nguyên, $\lambda$ là bước sóng)?
\choice
{$d_2 - d_1=\dfrac{k\lambda}{2}$}
{\True $d_2 - d_1=k\lambda$ }
{$d_2 - d_1=2k\lambda$ }
{$d_2 - d_1=(k + \dfrac{1}{2}) \lambda$ }
\loigiai{
Hai điểm dao động với biên độ cực đại: $\dfrac{2\pi(d_1-d_2)}{\lambda} = k2\pi\Rightarrow d_2-d_1 = k\lambda$.
}
\end{ex} \haidong{20}
\begin{ex}%[3P2Y5-1][Loai1] 
Trên mặt nước có hai nguồn kết hợp dao động điều hòa cùng pha theo phương thẳng đứng. Coi biên độ sóng không đổi khi sóng truyền đi. Trên mặt nước, trong vùng giao thoa, phần tử tại M dao động với biên độ cực đại khi hiệu đường đi của hai sóng từ hai nguồn truyền tới M bằng
\choice
{một số nguyên lần nửa bước sóng}
{một số lẻ lần nửa bước sóng}
{\True một số nguyên lần bước sóng}
{một số lẻ lần một phần tư bước sóng}
\loigiai{
Với hai nguồn cùng pha, phần tử tại M dao động với biên độ cực đại khi hiệu đường đi của hai sóng từ hai nguồn truyền tới M bằng một số nguyên lần bước sóng.
}
\end{ex} \haidong{20}
\begin{ex}
Khi hai nguồn sóng kết hợp giao thoa trên mặt nước, tại một điểm quan sát thấy biên độ dao động cực đại. Điều nào sau đây chắc chắn đúng tại điểm đó?
\choice
{\True Hiệu đường đi từ hai nguồn đến điểm đó bằng một số nguyên lần bước sóng.}
{Hai sóng đến điểm đó ngược pha nhau.}
{Khoảng cách từ hai nguồn đến điểm đó bằng nhau.}
{Hiệu đường đi từ hai nguồn đến điểm đó bằng một số lẻ bội số của nửa bước sóng.}
\loigiai{
Tại cực đại giao thoa, điều kiện là hiệu đường đi $\Delta d = d_2 - d_1 = k\lambda$ (với $k=0,\pm1,\pm2,\ldots$). Khi đó hai sóng đến điểm M cùng pha hoặc lệch pha $2\pi k$ nên biên độ tăng cường nhau.
}
\end{ex}

\begin{ex}
Hai sóng phát ra từ hai nguồn kết hợp cùng pha. Cực đại giao thoa nằm tại các điểm có
hiệu khoảng cách tới hai nguồn sóng bằng
\choice
{một ước số của bước sóng}
{\True một bội số nguyên của bước sóng}
{một bội số lẻ của nửa bước sóng}
{một ước số của nửa bước sóng}
\loigiai{}
\end{ex}
\begin{ex}
Trong thí nghiệm giao thoa sóng nước. Hai điểm trên cùng một phương truyền sóng, cách nhau một khoảng bằng bước sóng có dao động
\choice
{lệch pha $\dfrac{\pi}{2}$}
{ngược pha}
{lệch pha $\dfrac{\pi}{4}$}
{\True cùng pha}
\loigiai{}
\end{ex}
\loai{Vị trí cực tiểu}
\begin{ex}%[3P2Y5-1][Loai1] 
Giao thoa sóng ở mặt nước với hai nguồn kết hợp dao động điều hòa cùng pha theo phương thẳng đứng. Sóng truyền trên mặt nước có bước sóng $\lambda$. Cực tiểu giao thoa nằm tại những điểm có hiệu đường đi của hai sóng từ hai nguồn tới đó thỏa mãn điều kiện nào sau đây?
\choice
{$d_2-d_1=\left( k+\dfrac{1}{2} \right)\dfrac{\lambda}{2}$ với $k=0;\pm 1;\pm 2;...$}
{\True $d_2-d_1=\left( k+\dfrac{1}{2} \right)\lambda $ với $k=0;\pm 1;\pm 2;...$}
{$d_2-d_1=k\lambda $ với $k=0;\pm 1;\pm 2;...$}
{$d_2-d_1=k\dfrac{\lambda}{2}$ với $k=0;\pm 1;\pm 2;...$}
\loigiai{
Trong hiện tượng giao thoa của hai nguồn cùng pha, cực tiểu giao thoa có hiệu khoảng cách đến hai nguồn thỏa mãn $d_2-d_1=\left( k+\dfrac{1}{2} \right)\lambda $.
}
\end{ex} \haidong{20}
\begin{ex}%[3P2B5-1][Loai1]
\nguon{QG - 2017} Giao thoa ở mặt nước với hai nguồn sóng kết hợp đặt tại A và B dao động điều hòa cùng pha theo phương thẳng đứng. Sóng truyền tới mặt nước có bước sóng $\lambda$ . Cực tiểu giao thoa nằm tại những điểm có hiệu đường đi của hai sóng từ hai nguồn đến điểm đó bằng
\choice
{$(2k+1)\lambda $ với $k=0,\pm 1,\pm 2..$}
{\True $(k+0,5)\lambda $ với $k=0,\pm 1,\pm 2..$}
{$k\lambda $ với $k=0,\pm 1,\pm 2..$}
{$2k\lambda $ với $k=0,\pm 1,\pm 2..$}
\loigiai{
Các vị trí cực tiểu có hiệu khoảng cách bằng một số bán nguyên lần bước sóng $\Delta d=(k+0,5)\lambda$.
}
\end{ex} \haidong{20}
\begin{ex}%[3P2Y5-1][Loai1]
Ở mặt nước có hai nguồn sóng dao động theo phương vuông góc với mặt nước, có cùng phương trình $u = A\cos\omega t$. Trong miền gặp nhau của hai sóng, những điểm mà ở đó các phần tử nước dao động với biên độ cực tiểu sẽ có hiệu đường đi của sóng từ hai nguồn đến đó bằng
\choice
{\True một số lẻ lần nửa bước sóng}
{một số nguyên lần bước sóng}
{một số nguyên lần nửa bước sóng}
{một số lẻ lần bước sóng}
\loigiai{
Hai nguồn cùng pha, những điểm mà ở đó các phần tử nước dao động với biên độ cực tiểu sẽ có hiệu đường đi của sóng từ hai nguồn đến đó bằng một số lẻ lần nửa bước sóng.
}
\end{ex} \haidong{20}
\begin{ex}
Trong hiện tượng giao thoa sóng trên mặt nước, khoảng cách giữa hai cực đại liên tiếp nằm trên đường nối hai tâm sóng
\choice
{bằng hai lần bước sóng}
{bằng một bước sóng}
{\True bằng một nửa bước sóng}
{bằng một phần tư bước sóng}
\loigiai{}
\end{ex}
\loai{Loại khác}
\begin{ex}%[3P2Y5-1][Loai1] 
Để khảo sát giao thoa sóng cơ, người ta bố trí trên mặt nước nằm ngang hai nguồn kết hợp $S_1$ và $S_2$. Hai nguồn này dao động điều hòa theo phương thẳng đứng. Xem biên độ sóng không thay đổi trong quá trình truyền sóng. Tại trung điểm của đoạn $S_1S_2$, phần tử nước dao động với biên độ cực đại. Hai nguồn sóng đó dao động
\choice
{\True cùng pha nhau}
{lệch pha nhau góc $\dfrac{\pi}{3}$}
{ngược pha nhau}
{lệch pha nhau góc $0,5\pi$}
\loigiai{
\begin{itemize} 
\item Tại trung điểm của hai nguồn, ta có hiệu đường đi đến hai nguồn $\Delta d=0$.
\item Từ điều kiện để có cực đại giao thoa với hai nguồn kết hợp cùng pha: $\text{d}_1-d_2=k\lambda $.
\item Với $k=0$ ta thu được $\Delta \text{d}=0$.
\end{itemize}
}
\end{ex} \haidong{20}
\begin{ex}%[3P2Y5-1][Loai1]
Để khảo sát giao thoa sóng cơ, người ta bố trí trên mặt nước nằm ngang hai nguồn kết hợp $S_1$ và $S_2$. Hai nguồn này dao động điều hòa theo phương thẳng đứng, cùng pha. Xem biên độ sóng không thay đổi trong quá trình truyền sóng. Các điểm thuộc mặt nước và nằm trên đường trung trực của đoạn $S_1S_2$ sẽ
\choice
{dao động với biên độ bằng nửa biên độ cực đại}
{dao động với biên độ cực tiểu}
{\True dao động với biên độ cực đại}
{không dao động}
\loigiai{
Hai nguồn cùng pha, các điểm thuộc mặt nước và nằm trên đường trung trực của đoạn $S_1S_2$ sẽ dao động với biên độ cực đại.
}
\end{ex} \haidong{20}

\begin{ex}%[3P2Y5-1][Loai1] 
Giao thoa
\choice
{chỉ xảy ra khi ta thực hiện với sóng cơ}
{chỉ xảy ra khi ta thực hiện thí nghiệm trên mặt nước}
{\True là hiện tượng đặc trưng cho sóng}
{là sự chồng chất hai sóng trong không gian}
\loigiai{
Giao thoa là hiện tượng đặc trưng của sóng, xảy ra với cả sóng cơ và sóng điện từ.
}
\end{ex} \haidong{20}
\begin{ex}
Phát biểu nào sau đây là không đúng?
\choice
{Khi xảy ra hiện tượng giao thoa sóng trên mặt chất lỏng, tồn tại các điểm dao động với biên độ cực đại}
{Khi xảy ra hiện tượng giao thoa sóng trên mặt chất lỏng, tồn tại các điểm không dao động}
{Khi xảy ra hiện tượng giao thoa sóng trên mặt chất lỏng, các điểm không dao động tạo thành các vân cực tiểu}
{\True Khi xảy ra hiện tượng giao thoa sóng trên mặt chất lỏng, các điểm dao động mạnh tạo thành các đường thẳng cực đại}
\loigiai{}
\end{ex}
\begin{ex}
Khi xảy ra hiện tượng giao thoa sóng nước với hai nguồn kết hợp cùng pha A, B. Những điểm trên mặt nước nằm trên đường trung trực của AB sẽ
\choice
{\True dao động với biên độ lớn nhất}
{dao động với biên độ bé nhất}
{đứng yên không dao động}
{dao động với biên độ có giá trị trung bình}
\loigiai{}
\end{ex}

\PhanII
\setcounter{ex}{0}
\begin{ex}
Trên mặt nước, hai nguồn sóng kết hợp A, B cùng biên độ, cùng tần số, ngược pha nhau.
\choiceTF[t]
{\True Trung điểm của đoạn thẳng nối hai nguồn là một cực tiểu giao thoa}
{\True Đường thẳng trung trực của đoạn thẳng nối hai nguồn là một đường cực tiểu giao thoa}
{Trên đường thẳng nối hai nguồn, khoảng cách giữa một cực đại và cực tiểu liên tiếp là $\dfrac{\lambda}{2}$}
{Trên đường thẳng nối hai nguồn, khoảng cách giữa hai cực đại liên tiếp là $\lambda$}
\loigiai{
\begin{itemchoice}
\itemch
\itemch
\itemch
\itemch
\end{itemchoice}
}
\end{ex} \haidong{20}

\begin{ex}
Trên mặt nước, hai nguồn đồng bộ đang phát ra sóng.
\choiceTF[t]
{\True Các đường cực tiểu giao thoa luôn là các hyperbol}
{Các điểm trên mặt nước thuộc đường trung trực của đoạn thẳng nối 2 nguồn thì dao động với biên độ cực đại}
{\True Các đường cực đại giao thoa cùng bậc đối xứng nhau qua đường trung trực đoạn thẳng nối hai nguồn}
{\True Trên các đường cực tiểu, sóng đang dao động với biên độ bé nhất}
\loigiai{
\begin{itemchoice}
\itemch
\itemch
\itemch
\itemch
\end{itemchoice}
}
\end{ex} \haidong{20}
\begin{ex}
Thực hiện giao thoa sóng cơ với hai nguồn phát sóng đồng bộ.
\choiceTF[t]
{\True Trên đoạn thẳng nối hai nguồn, số cực đại luôn chênh lệch số cực tiểu một điểm}
{Số điểm không dao động trên đoạn nối hai nguồn là số lẻ}
{Số điểm dao động mạnh nhất trên đoạn nối hai nguồn là số chẵn}
{\True Các vân cực đại và các vân cực tiểu nằm xen kẽ nhau trên trường giao thoa}
\loigiai{
\begin{itemchoice}
\itemch
\itemch
\itemch
\itemch
\end{itemchoice}
}
\end{ex} \haidong{20}
\begin{ex}
Giao thoa sóng trên mặt nước của hai nguồn kết hợp có cùng biên độ và cùng pha với nhau.
\choiceTF[t]
{Trung điểm của đoạn thẳng nối hai nguồn là một cực tiểu giao thoa}
{Đường trung trực nối hai nguồn là tập hợp các cực đại giao thoa bậc một}
{\True Tập hợp các cực tiểu giao thoa cùng bậc luôn tạo thành một đường hyperbol}
{Số cực đại giao thoa trên đoạn thẳng nối hai nguồn luôn là một số chẵn}
\loigiai{
\begin{itemchoice}
\itemch
\itemch
\itemch
\itemch
\end{itemchoice}
}
\end{ex} \haidong{20}